\chapter{Static D2NN Fundamental Limit}

\section{Problem Framing}

This chapter answers the field-level question left open by the optics crisis.
Suppose a receiver-side phase-only D2NN is modeled as a sampled, scalar,
monochromatic, lossless operator $H$. Can such a fixed operator improve the
field similarity between a turbulent realization $U_t$ and its vacuum reference
$U_v$?

The scope is intentionally narrow. The operator $H$ in this chapter refers only
to the sampled phase-mask plus band-limited angular spectrum propagation block
on the propagating subspace. Zero-padding, finite detector windows, cropping,
and focal-lens measurements are outside the theorem statements below.

\section{Unitary Block Model}

Each layer is a phase-only multiplier and each inter-layer propagation step is a
band-limited angular spectrum propagator. Under the model used in this project,
their composition is unitary on the sampled propagating subspace.

\begin{definition}[Complex overlap used in Theorem 1]
For two sampled complex fields $U_1$ and $U_2$ on the same receiver grid, the
complex overlap is
\begin{equation}
\CO(U_1,U_2)=
\frac{|\langle U_1,U_2\rangle|}{\|U_1\|_2\,\|U_2\|_2}.
\end{equation}
Mixed quantities such as $\CO(HU_t,U_v)$ are different observables and are not
covered by the invariance theorem.
\end{definition}

\begin{theorem}[Complex-overlap invariance]
Let $H$ be a unitary phase-only D2NN block. Then
\begin{equation}
\CO(HU_t, HU_v) = \CO(U_t, U_v).
\end{equation}
\end{theorem}

\begin{proof}
The complex overlap is the normalized inner product. Because $H^\dagger H = I$,
\begin{equation}
\langle HU_t, HU_v \rangle = \langle U_t, H^\dagger H U_v \rangle
= \langle U_t, U_v \rangle.
\end{equation}
Norms are also preserved, so the normalized ratio is unchanged.
\end{proof}

\begin{theorem}[Distance invariance]
If $H$ is unitary, then
\begin{equation}
\|HU_t - HU_v\|_2 = \|U_t - U_v\|_2.
\end{equation}
\end{theorem}

\begin{proof}
Expand the squared norm and use $H^\dagger H = I$:
\begin{align}
\|HU_t - HU_v\|_2^2
&= \langle H(U_t-U_v), H(U_t-U_v) \rangle \\
&= \langle U_t-U_v, H^\dagger H (U_t-U_v) \rangle \\
&= \|U_t-U_v\|_2^2.
\end{align}
\end{proof}

\section{What Is Not Invariant}

Intensity observables are different. A detector-plane bucket metric can be
written as
\begin{equation}
\PIB(U) = \frac{\langle U, B U \rangle}{\langle U, U \rangle},
\end{equation}
where $B$ is a quadratic observable selecting the bucket support. After a
unitary transform,
\begin{equation}
\PIB(HU) = \frac{\langle U, H^\dagger B H U \rangle}{\langle U, U \rangle},
\end{equation}
which only equals the original value when $H^\dagger B H = B$. In general, that
condition does not hold. This is why a static D2NN can alter detector-plane
intensity metrics without improving field-level similarity.

\section{Numerical Verification}

\runfigure{0327-theorem-verify-defocus-1layer/deterministic_verification.png}
  {Deterministic defocus verification}
  {A deterministic defocus aberration can be corrected by a simple phase-only
  optical element. This confirms that the later random-turbulence limitation is
  not a statement about every deterministic distortion.}

\runfigure{0329-paper-figures-static-d2nn/theorem1_co_verification.png}
  {Theorem 1 verification}
  {Numerical verification that $\CO(U_t,U_v)$ and $\CO(HU_t,HU_v)$ are
  preserved to numerical precision, while mixed comparisons such as
  $\CO(HU_t,U_v)$ can change dramatically.}

\runfigure{0329-paper-figures-static-d2nn/theorem2_l2_verification.png}
  {Theorem 2 verification}
  {Numerical verification of $L_2$-distance preservation for the same
  strategies.}

\runfigure{0329-paper-figures-static-d2nn/theorem_verification_table.png}
  {Invariant summary}
  {Summary of preserved versus non-preserved quantities in the sampled unitary
  model.}

\runfigure{0329-paper-figures-static-d2nn/fig3_deterministic_vs_random.png}
  {Deterministic versus random distortion}
  {Deterministic aberrations can be conjugated, whereas random turbulence still
  respects the unitary invariants above.}

\begin{table}[H]
  \centering
  \begin{tabular}{lccc}
    \toprule
    \smalltablehead{Strategy} & \smalltablehead{$\CO(HU_t,HU_v)$} & \smalltablehead{Mixed $\CO(HU_t,U_v)$} & \smalltablehead{$L_2$ verdict} \\
    \midrule
    PIB only & 0.304446 & 0.014716 & preserved \\
    CO+PIB hybrid & 0.304446 & 0.190252 & preserved \\
    Strehl only & 0.304446 & 0.023102 & preserved \\
    Intensity overlap & 0.304446 & 0.183928 & preserved \\
    \bottomrule
  \end{tabular}
  \caption{Field-level invariants are preserved, while mixed comparisons can
  vary widely. The preserved value $0.304446$ is taken from
  \texttt{theorem\_verification\_results.json}.}
\end{table}

\section{Interpretation}

\begin{discovery}
Within the stated sampled unitary model and for field-level similarity, a
static phase-only D2NN does not act as a random-turbulence wavefront corrector.
It is better interpreted as a passive mode converter or detector-shaping optical
front end.
\end{discovery}

This conclusion does \emph{not} say that every intensity-based observable is
fixed. On the contrary, detector-plane quantities such as PIB can change
substantially. The mistake would be to interpret that change as direct evidence
of field recovery. That exact mistake is what the next chapter diagnoses in the
legacy focal-PIB sweep.
